% =====================================================================
%  theme.tex — tiskarski sloj teme (xelatex / scrbook)
%  Uključuje se preko _quarto-pdf.yml: include-in-header.
% =====================================================================
%  Preslikava design-tokens.yml i styles/_tokens.scss. Provjera sinkronizacije:
%      Rscript scripts/check-tokens.R
%
%  PLACEHOLDER: paleta je neutralna, a pisma su LaTeX-ova zadana (Latin
%  Modern) da PDF gradi bez ijedne vanjske datoteke. Blok za brendirana
%  pisma je pripremljen i zaključan dolje. Vidi DESIGN.md.
% =====================================================================

% ---- paleta (JEDINO mjesto s bojama u tiskarskom sloju) -------------
\definecolor{paper}{HTML}{FFFFFF}       % token:paper
\definecolor{papersoft}{HTML}{F4F4F5}   % token:paper-soft
\definecolor{ink}{HTML}{18181B}         % token:ink
\definecolor{inksoft}{HTML}{3F3F46}     % token:ink-soft
\definecolor{inkmute}{HTML}{71717A}     % token:ink-mute
\definecolor{inkfaint}{HTML}{A1A1AA}    % token:ink-faint
\definecolor{accent}{HTML}{2563A8}      % token:accent
\definecolor{accentdeep}{HTML}{1B4A7D}  % token:accent-deep
\definecolor{accentwash}{HTML}{E8F0F8}  % token:accent-wash
\definecolor{emphasis}{HTML}{9C2B2B}    % token:emphasis
\definecolor{archival}{HTML}{8A6D3B}    % token:archival

% ---- boja papira i teksta -------------------------------------------
\pagecolor{paper}
\AtBeginDocument{\color{ink}}
\definecolor{shadecolor}{HTML}{F4F4F5}   % podloga blokova koda -> paper-soft

% ---- pisma ----------------------------------------------------------
\usepackage{fontspec}
% PLACEHOLDER: bez \setmainfont, pa xelatex koristi Latin Modern i PDF se
% gradi na svakom stroju. Kad tipografija bude odabrana, stavite statičke
% instance u fonts/print/ i otključajte blok ispod (uskladite s design-tokens.yml).
%
% \setmainfont{NazivSerifa}[
%   Path = fonts/print/ , Extension = .ttf ,
%   UprightFont = *-Regular , ItalicFont = *-Italic ,
%   BoldFont = *-SemiBold , BoldItalicFont = *-BoldItalic ]
% \setsansfont{NazivSansa}[
%   Path = fonts/print/ , Extension = .ttf ,
%   UprightFont = *-Regular , ItalicFont = *-Italic , BoldFont = *-Bold ]
% \setmonofont{NazivMona}[
%   Path = fonts/print/ , Extension = .ttf ,
%   UprightFont = *-Regular , BoldFont = *-SemiBold , Scale = 0.86 ]

% Natpis u velikim slovima s razmakom (eyebrow). Bez fontspec LetterSpace
% opcije, pa radi i sa zadanim pismom.
\newcommand{\eyebrow}[2][accent]{{\sffamily\bfseries\footnotesize\color{#1}\MakeUppercase{#2}}}

% ---- poveznice ------------------------------------------------------
\AtBeginDocument{\hypersetup{
  linkcolor=accentdeep, citecolor=accentdeep,
  urlcolor=accentdeep, filecolor=accentdeep }}

% =====================================================================
%  NASLOVI
% =====================================================================

% ---- DIO (part) -----------------------------------------------------
\setkomafont{part}{\normalfont\color{ink}}
\renewcommand*{\partformat}{}
\renewcommand*{\partheadstartvskip}{\null\vskip 8em}
\renewcommand*{\partheadendvskip}{\vfil\newpage}
\newcommand{\BookPartHead}[1]{%
  \raggedright
  {\color{archival}\rule{\linewidth}{0.8pt}}\par
  \vskip 0.85em
  {\bfseries\color{ink}\fontsize{28}{33}\selectfont #1\par}%
  \vskip 0.85em
  {\color{accent}\rule{\linewidth}{1.5pt}}%
}
\renewcommand{\partlineswithprefixformat}[3]{\BookPartHead{#3}}

% ---- POGLAVLJE ------------------------------------------------------
% Natpis "Poglavlje" gasi se u dodacima da literatura ne bude krivo označena.
\newif\ifbookeyebrow \bookeyebrowtrue
\let\bookoldappendix\appendix
\renewcommand{\appendix}{\bookoldappendix\bookeyebrowfalse}
\setkomafont{chapter}{\bfseries\color{ink}}
\renewcommand*{\chapterformat}{}
\renewcommand*{\chapterheadstartvskip}{\null\vskip 5.5em}
\newcommand{\BookChapHead}[1]{%
  \raggedright
  \ifbookeyebrow {\eyebrow{Poglavlje}\par}\vskip 0.55em \fi
  {\usekomafont{chapter}\fontsize{22}{26}\selectfont #1\par}%
  \vskip 0.45em
  {\color{accent}\rule{\linewidth}{1.3pt}}%
  \par\vskip 2.2em
}
\renewcommand{\chapterlinesformat}[3]{\BookChapHead{#3}}
\renewcommand{\chapterlineswithprefixformat}[3]{\BookChapHead{#3}}

% ---- ODJELJCI -------------------------------------------------------
\setkomafont{section}{\bfseries\color{accentdeep}}
\setkomafont{subsection}{\bfseries\color{accentdeep}}
\setkomafont{subsubsection}{\sffamily\bfseries\color{inksoft}}
\setkomafont{paragraph}{\sffamily\bfseries\color{inksoft}}

% ---- sadržaj --------------------------------------------------------
\setkomafont{partentry}{\bfseries\color{ink}}
\setkomafont{chapterentry}{\color{ink}}

% =====================================================================
%  ŽIVA PAGINA
% =====================================================================
\usepackage[automark]{scrlayer-scrpage}
\clearpairofpagestyles
\automark[section]{chapter}
\setkomafont{pagehead}{\sffamily\footnotesize\color{inkmute}}
\setkomafont{pagenumber}{\normalfont\ttfamily\footnotesize\color{inkmute}}
\setkomafont{headsepline}{\color{accent}}
\lehead{{\normalfont\sffamily\upshape\footnotesize\color{inkmute}\leftmark}}
\rohead{{\normalfont\sffamily\upshape\footnotesize\color{inkmute}\rightmark}}
\cfoot*{{\ttfamily\footnotesize\color{inkmute}\thepage}}
\KOMAoptions{headsepline=0.6pt}
\pagestyle{scrheadings}

% =====================================================================
%  NUMERACIJA TEOREMA — ravna i kontinuirana
%  Odjeljci su nenumerirani, pa je zadano "Definicija 0.0.1" besmisleno.
% =====================================================================
\usepackage{remreset}
\makeatletter
\newcommand{\bk@flatten}[1]{%
  \@ifundefined{c@#1}{}{\@removefromreset{#1}{section}\@namedef{the#1}{\arabic{#1}}}}
\AtBeginDocument{%
  \bk@flatten{definition}\bk@flatten{proposition}\bk@flatten{theorem}%
  \bk@flatten{lemma}\bk@flatten{corollary}\bk@flatten{example}\bk@flatten{exercise}%
}
\makeatother

% =====================================================================
%  PEDAGOŠKE KUTIJE  (omata ih pdf-filters/book-callouts.lua)
% =====================================================================
\usepackage[most]{tcolorbox}
\newtcolorbox{bookcallout}[2]{%
  enhanced, breakable, sharp corners,
  colback=papersoft, colframe=#1,
  boxrule=0.4pt, leftrule=3pt,
  left=14pt, right=13pt, top=11pt, bottom=12pt,
  before skip=12pt, after skip=12pt,
  fontupper=\color{ink},
  coltitle=#1, fonttitle=\sffamily\bfseries\footnotesize,
  attach title to upper={\par\vskip 7pt},
  title={#2}%
}
% Natpisi se predaju VEĆ VELIKIM SLOVIMA, bez \MakeUppercase. Razlog je Đ:
% automatsko povećavanje slova ispisalo ga je kao goli D („NADITE GREŠKU"),
% dok literalno Đ prolazi ispravno.
\newenvironment{calloutvinjeta}{\begin{bookcallout}{archival}{VINJETA}}{\end{bookcallout}}
\newenvironment{calloutdivljina}{\begin{bookcallout}{emphasis}{STATISTIKA U DIVLJINI}}{\end{bookcallout}}
\newenvironment{calloutmodel}{\begin{bookcallout}{accent}{PITAJTE MODEL}}{\end{bookcallout}}
\newenvironment{calloutgreska}{\begin{bookcallout}{emphasis}{NAĐITE GREŠKU}}{\end{bookcallout}}

% =====================================================================
%  OPISI SLIKA I TABLICA
% =====================================================================
\setkomafont{caption}{\sffamily\small\color{inkmute}}
\setkomafont{captionlabel}{\sffamily\bfseries\color{accentdeep}}

% =====================================================================
%  NASLOVNICA (nadjačava \maketitle; ide prije sadržaja)
% =====================================================================
\makeatletter
\renewcommand{\maketitle}{%
  \begin{titlepage}%
  \newgeometry{margin=2.6cm}%
  \raggedright
  \null\vskip 4.5em
  \eyebrow[inkmute]{Udžbenik statistike}\par
  \vskip 0.9em
  {\color{accent}\rule{\linewidth}{1pt}}\par
  \vskip 1.6em
  {\bfseries\color{ink}\fontsize{34}{40}\selectfont \@title\par}%
  \ifx\@subtitle\@empty\else
    \vskip 0.8em
    {\itshape\fontsize{15}{20}\selectfont\color{inksoft}\@subtitle\par}%
  \fi
  \vfill
  {\color{archival}\rule{\linewidth}{0.6pt}}\par
  \vskip 0.9em
  {\sffamily\large\color{ink}%
    \renewcommand{\and}{\ \textcolor{archival}{\textperiodcentered}\ }\@author\par}%
  \vskip 0.5em
  {\rmfamily\color{inkmute}\@date\par}%
  \restoregeometry
  \end{titlepage}%
}
\makeatother
